\section{Kết luận}
\label{sec:conclusion}

Nghiên cứu này trình bày một pipeline CausalRAG cho hỏi--đáp pháp luật hình
sự Việt Nam, trong đó đối tượng truy hồi trung tâm là đồ thị điều kiện--hệ
quả pháp luật và đối tượng xác minh là một mô hình quy tắc có thể thực thi.
Pipeline kết hợp chuẩn hóa quy tắc, causal memory, truy hồi event/rule nhiều
bước, xếp hạng causal path, query-aware verification và sinh câu trả lời
extractive có citation.

LegalSCM được xây dựng như một mô hình nhân quả cấu trúc pháp lý tất định ba
giá trị. Mô hình có thể xác nhận factual path, thực hiện hard intervention
trên mediator và suy luận lại outcome trong phạm vi graph đã mã hóa. Song
song với đó, path ablation được giữ như một baseline reachability sau khi
xóa node và không được xem là cách cài đặt phép $do(\cdot)$.

Trên benchmark $250$ câu hỏi, pipeline đạt Rule Recall@5 $69{,}93\%$, Event
Recall@5 $74{,}66\%$, Top-1 Exact Path $35{,}20\%$, Oracle Exact Path
$71{,}20\%$, Verification Accuracy $92{,}00\%$, Answer Token F1
$68{,}69\%$, ROUGE-L F1 $67{,}67\%$ và Citation F1 $62{,}86\%$. Khoảng
cách $36$ điểm phần trăm giữa Top-1 và Oracle Exact Path cho thấy path
ranking là điểm nghẽn chính. Reverse reasoning là nhóm yếu nhất, trong khi
metric path tuyến tính chưa phản ánh phù hợp branch và convergence.

Paired ablation không ghi nhận bất kỳ chênh lệch quality metric nào giữa
structural SCM và path ablation. Đây là một null quality result, không phải
bằng chứng về ưu thế của structural verification hoặc sự tương đương lý
thuyết giữa hai phương pháp. Structural SCM cung cấp intervention semantics
và world-state trace chi tiết hơn, nhưng cần thời gian trung bình trên mỗi
mẫu cao hơn khoảng $3{,}23$ lần. Benchmark hiện tại, đặc biệt là $20$ mẫu
direct-edge negative, chưa đủ nhạy để định lượng lợi ích accuracy của hard
intervention.

Đóng góp chính của nghiên cứu vì vậy nằm ở việc xây dựng và kiểm tra một
pipeline có ranh giới tuyên bố rõ ràng, cung cấp paired artifact có thể truy
vết và xác định chính xác các nút thắt còn lại. Bước tiếp theo có tính quyết
định không phải chỉ là tăng một aggregate score, mà là xây dựng benchmark
structural-counterfactual cân bằng, được chuyên gia thẩm định, đồng thời cải
thiện query-aware path ranking và reverse retrieval. Cho đến khi hoàn thành
các bước đó, hệ thống cần được xem là nguyên mẫu nghiên cứu hỗ trợ phân tích,
không phải công cụ tư vấn hoặc ra quyết định pháp lý.